\scanpage{115}
\mspage{102}
% archive sheet number: 114

\inserted{bien} une catégorie de Quillen spéciale close, est
$\pm$ tautologique. ].

\uline{Prouvons} b). La relation (6.1.6) est
tautologique. Les relations (6.1.7) sont
contenues ds 5.8. (cf $2^{\unsure{\mathrm{me}|\mathrm{e}}}$ relation
ds (5.8.3) pour $C'=\unsure{v|u}_{!}(C)^{\sim}$ ou
$C(I)=(v_{!}(C(I_{0})))^{\sim}$, et (5.8.5) pour
$TC(I)=(v_{!}(TC(I_{0})))^{\sim}$.

\uline{Prouvons} c). La première assertion
est, par Quillen, un cas particulier
des relations (6.1.10) à (6.1.14), c'est
elles qu'il s'agit d'établir. Les relations
(6.1.13) et par suite (6.1.14) sont
tautologiques en vertu des définitions,
\deleted{\unsure{donc}} il s'agit d'établir \uline{(6.1.10)}
et (6.1.12).

Pour les formules (6.1.10), j'utilise
le diagramme commutatif
\diagram{\begin{tikzcd}[ampersand replacement=\&,column sep=large]
I \arrow[d,dash,"{u}"'] \& I_{0} \arrow[l,"{v}"'] \arrow[d,"{u_{0}}"]\\
I' \& I'_{0} \arrow[l,"{v'}"']
\end{tikzcd}},
et je note que les foncteurs $u_{!}^{M},u_{0!}^{M}$, commutant aux petites limites inductives, \deleted{\unsure{comme}} \inserted{se}

